% dissemination_body.tex

% ============================================================
\section*{Dissemination Strategy Guide}
% ============================================================

\subsection*{Project Information}

\begin{tabular}{@{}P{3.5cm}P{11cm}}
\textbf{Title:}       & \#TERRA. Paths to Sustainability in Adolescence \\[4pt]
\textbf{Reference:}   & COMPETE2030-FEDER-00903800 \\[4pt]
\textbf{Funding:}     & COMPETE 2030 and Portugal 2030 programmes; European Union \\[4pt]
\textbf{Host:}        & Faculty of Psychology and Education Sciences at the University of Porto (FPCEUP); Center for Psychology at the University of Porto (CPUP) \\[4pt]
\textbf{Start date:}  & September 2025 \\[4pt]
\textbf{End date:}    & September 2028 \\[4pt]
\end{tabular}

\vspace{0.5cm}

\noindent\textbf{Research team:}
\begin{itemize}
  \item Tiago Ferreira (PI) | \href{mailto:tiagoferreira@fpce.up.pt}{tiagoferreira@fpce.up.pt}
  \item Filipa Nunes | \href{mailto:fnunes@fpce.up.pt}{fnunes@fpce.up.pt}
  \item H\'{e}lder Sp\'{i}nola | \href{mailto:hspinola@staff.uma.pt}{hspinola@staff.uma.pt}
  \item Marisa Matias | \href{mailto:marisa@fpce.up.pt}{marisa@fpce.up.pt}
  \item Paula Mena Matos | \href{mailto:pmmatos@fpce.up.pt}{pmmatos@fpce.up.pt}
\end{itemize}

\vspace{0.5cm}

\noindent\textit{This Dissemination Strategy Guide is part of project \#TERRA. Paths to Sustainability in Adolescence, which was funded by the COMPETE 2030 and Portugal 2030 programmes and co-financed by the European Union (COMPETE2030-FEDER-00903800).}

\vspace{0.8cm}

% ============================================================
\subsection*{Short Background}
% ============================================================

\subsubsection*{The project}

The \#TERRA project examines pro-environmental attitudes and behaviors during adolescence, with a focus on the individual and contextual factors that shape them. It combines a nationwide cross-sectional study with a longitudinal approach to track developmental changes and daily variations in ecological behavior. The main goals are:

\begin{itemize}
  \item Develop and validate tools to measure these attitudes and behaviors in areas like sustainable consumption, waste management, and energy efficiency.
  \item Identify the main factors that shape adolescents' pro-environmental attitudes and behaviors.
  \item Examine how these attitudes and behaviors evolve during adolescence and what predicts change.
  \item Analyze daily fluctuations considering personal factors (e.g., mood, life satisfaction) and immediate social and environmental contexts.
\end{itemize}

\subsubsection*{This document}

This dissemination strategy is a guide for action that presents the strategy for communication and dissemination of \#TERRA. First, we present the project identity (logo) to establish the basic principles to preserve visual properties and facilitate the correct dissemination, perception, identification, and memorization of the \#TERRA communication material. Next, we present the goals of the communication strategy and the dissemination tools of \#TERRA.

% ============================================================
\section{Introduction}
% ============================================================

\subsection{About the Project}

The \#TERRA project aims to investigate Pro-Environmental (PE) attitudes and behaviors during adolescence, seeking to identify their individual and contextual antecedents. This project combines a cross-sectional analysis of adolescents' PE attitudes and behaviors, based on a representative sample of the Portuguese population, with a longitudinal analysis focused on individual changes associated with age and daily fluctuations in ecological behaviors. This project aims to:

\begin{enumerate}[label=(\arabic*)]
  \item develop, adapt, and validate assessment instruments that reliably and validly measure PE attitudes and behaviors in adolescence, encompassing both the perceptual and general dimensions of actions, as well as their frequency and nature within ecologically relevant domains such as sustainable consumption, waste management, and energy efficiency;
  \item characterize and identify the main individual and contextual factors that shape the PE attitudes and behaviors of Portuguese adolescents;
  \item examine the developmental trajectory of PE attitudes and behaviors throughout adolescence, identifying predictive and risk factors associated with interindividual changes;
  \item analyze daily and weekly fluctuations in adolescents' ecologically relevant behaviors, investigating the role of intraindividual predictors (e.g., mood states and life satisfaction) and immediate contextual factors (e.g., daily social and environmental context) that shape variations in these behaviors.
\end{enumerate}

Taken together, research results will enrich research in basic psychology, through a deeper understanding of the complex interaction between individual and contextual predictors of PE attitudes and behaviors, providing critical evidence to promote environmentally sustainable behaviors among young people. The \#TERRA project will also have a significant social impact, with direct implications across multiple intervention domains. In the educational domain, it will support the design of educational programs aimed at promoting sustainable behaviors, contributing to more effective pedagogical practices in fostering environmentally responsible citizens. In the domain of public policy, the findings will provide solid scientific guidance for national climate change mitigation strategies, enabling more informed decisions aligned with the United Nations Sustainable Development Goals (SDGs). Furthermore, the \#TERRA project will impact social awareness by offering evidence to underpin campaigns that strengthen active environmental citizenship and encourage the adoption of more responsible consumption habits.

\subsection{Visual Identity: Logo}

The logo of the \#TERRA project combines two main elements:

\begin{itemize}
  \item The project name, ``\#TERRA'', is written in bold letters with a hashtag that emphasizes its social and digital dimension.
  \item A symbolic image of a water droplet enclosing a globe, representing the link between Earth and sustainability. The droplet suggests caring for natural resources, while the globe highlights the global scope of ecological challenges.
\end{itemize}

The logo is available in four versions:

\begin{itemize}
  \item \textbf{Green or gradient version} -- conveys nature, growth, and ecological responsibility.
  \item \textbf{Black-and-white version} -- ideal for high-contrast applications and minimalist contexts.
  \item \textbf{Gradient version with tagline} -- includes ``Paths to Sustainability in Adolescence,'' reinforcing the project's focus on youth and sustainability.
  \item \textbf{Black-and-white version with tagline} -- for simple, formal, or monochromatic materials.
\end{itemize}

\begin{figure}[h!]
  \centering
  
\includegraphics[width=0.22\textwidth]{Prancheta 2.png}\hfill
  
\includegraphics[width=0.22\textwidth]{Prancheta 3.png}\hfill
  \includegraphics[width=0.22\textwidth]{Prancheta 6.png}\hfill
  \includegraphics[width=0.22\textwidth]{Prancheta 7.png}
  \caption{The four versions of the \#TERRA project logo: green (left), black-and-white, green with tagline, and black-and-white with tagline (right).}
  \label{fig:logos}
\end{figure}

These clean and modern designs reflect the project's commitment to environmental awareness and global sustainability.

\subsection{Utilization}

The \#TERRA logo should be used for all official communications of the project, namely the website, social media, communication material, reports, and external communications. The logo can be scaled to publications, flyers, and other materials. However, using the logo in too small a format is not recommended, as it would no longer be recognizable.

\subsection{Funding Organization}

In accordance with the provisions outlined in the Rules for Participation, the \#TERRA research team undertakes to explicitly recognize COMPETE 2030 and PORTUGAL 2030 as the project's funding organizations, as well as the European Union as a co-funder of the project. The research team undertakes to communicate these programmes in all project-related communications, activities, documents, and materials, as well as at any promotional events relating to the project. This includes, without limitation, the annual report, website content, emails, social media posts, digital materials, and communications on social networks.

Please see below the project email footer:

\begin{figure}[h!]
  \centering
  \includegraphics[width=\textwidth]{emailfooter.png}
  \caption{Standard \#TERRA email footer including funding acknowledgement and programme logos.}
  \label{fig:emailfooter}
\end{figure}

The following image shows the footer of the project website:

\begin{figure}[h!]
  \centering
  \includegraphics[width=\textwidth]{websitefooter.png}
  \caption{Footer of the \#TERRA project website showing funder logos and acknowledgement statement.}
  \label{fig:websitefooter}
\end{figure}

% ============================================================
\section{Dissemination Strategy Guide}
% ============================================================

This Dissemination Strategy Guide outlines the approach for planning, targeting, developing, and implementing communication and dissemination actions throughout the \#TERRA project. Special attention will be given to promoting and sharing project results and ensuring accurate and timely information flow among the research team, advisory panel, stakeholders, and, most importantly, the broader community. The \#TERRA project views research as a tool for environmental awareness. Therefore, measures will be taken to strengthen the connection between science and society, fostering dialogue and engagement with adolescents, educators, policymakers, and the community. This plan defines the project's dissemination objectives and the actions required to achieve them. It establishes priorities and key goals, identifies main stakeholders, and specifies the most effective channels and tools for outreach. It also sets timelines for planned dissemination activities and stakeholder engagement, ensuring that communication efforts are strategic, inclusive, and impactful.

This Dissemination Strategy Guide covers both internal and external communication.

\subsection{Internal Communication}

Internal communication aims to maintain transparency, collaboration, and timely information exchange among team members and partners. The following mechanisms will be implemented:

\begin{itemize}
  \item \textbf{Email Communication:} All official internal communication will use the dedicated project email: \href{mailto:terra@fpce.up.pt}{terra@fpce.up.pt}.
  \item \textbf{Virtual and Presential Meetings:} Regular videoconferences and in-person meetings will be scheduled to discuss progress, address challenges, and plan upcoming activities.
  \item \textbf{Data Storage and Sharing:} Project data and documents will be securely stored and shared through the University of Porto Server and Center for Psychology at the University of Porto Network Attached Storage (NAS). These platforms ensure compliance with institutional data protection policies and facilitate collaborative work.
\end{itemize}

\subsection{External Communication}

External communication focuses on disseminating project results, promoting engagement, and ensuring visibility among diverse audiences, including adolescents, educators, policymakers, researchers, and the community. The following channels and actions will be implemented:

\begin{itemize}
  \item Project Website
  \item Project Social Media
  \item Open Science Framework (OSF)
  \item Participation in Scientific Events and Publications
  \item Infographics and Press Releases
  \item White Book and Closing Seminar
\end{itemize}

These channels and actions aim to guarantee that the project's outputs are accessible, transparent, and impactful.

\subsubsection{Project Website}

The first version of the project website is already available at: \href{https://projeto-terra.fpce.up.pt/en/home/}{projeto-terra.fpce.up.pt}. The website is bilingual (Portuguese and English), ensuring accessibility for both national and international audiences. It serves as the main hub for external communication and includes the following pages:

\begin{itemize}
  \item \textbf{Project:} A visually appealing banner featuring the project logo and tagline ``Paths to Sustainability in Adolescence'', reinforcing the project's identity and mission. The homepage also provides a summary of the project's objectives, allowing visitors to quickly understand its purpose and scope.
  \item \textbf{Team:} Short bios for each team member including institutional affiliations and ORCID, highlighting their expertise and role in the project. It also includes profiles of all Advisory Panel members, ensuring transparency and recognition of their contribution.
  \item \textbf{Announcements:} A dedicated page for announcements, events, and progress reports. This section will showcase key milestones, participation in conferences, and activities relevant to both academic and general audiences.
  \item \textbf{Stakeholders:} This page will present the strategic community partners involved in the project. It will highlight organizations and institutions that collaborate with \#TERRA to promote sustainability in adolescence, ensuring visibility and recognition of their role in supporting the project's objectives.
\end{itemize}

\begin{figure}[h!]
  \centering
  \includegraphics[width=\textwidth]{website.png}
  \caption{Screenshot of the \#TERRA project website homepage (\href{https://projeto-terra.fpce.up.pt/en/home/}{projeto-terra.fpce.up.pt}).}
  \label{fig:website}
\end{figure}

\subsubsection{Project Social Media}

Official profiles were created on the following platforms: \textbf{LinkedIn} and \textbf{Mastodon}. These platforms serve as key communication channels to promote project activities and results, as well as share relevant content on pro-environmentalism in adolescence (e.g., podcasts, forums, and news). The selection of these networks is based on their ability to reach diverse audiences and foster meaningful engagement:

\begin{itemize}
  \item \textbf{LinkedIn} (\href{https://www.linkedin.com/in/terra-project-1698403a3}{www.linkedin.com/in/terra-project-1698403a3}): A professional network ideal for connecting with educators, policymakers, and organizations interested in sustainability and youth development. It enhances the project's visibility among institutional stakeholders and facilitates knowledge exchange.
  \item \textbf{Mastodon} (\href{https://mastodon.social/@terra_project}{mastodon.social/@terra\_project}): A decentralized social network that promotes open dialogue and community engagement. It aligns with the project's values of inclusivity and sustainability, offering an alternative to mainstream platforms and fostering discussions on environmental issues.
\end{itemize}

\subsubsection{Open Science Framework (OSF)}

The OSF platform (\href{https://osf.io/utmak/overview}{osf.io/utmak/overview}) will host protocols, datasets, and materials, ensuring open access and compliance with transparency standards. This reinforces the project's commitment to Open Science principles.

\subsubsection{Participation in Scientific Events and Publications}

The project will actively engage in national and international scientific events to disseminate findings and foster dialogue with the academic community. Participation in conferences and symposia (at least 4 presentations) will provide opportunities to present research outcomes, share methodologies, and strengthen collaborations with experts in sustainability and adolescent development. In parallel, the project will prioritize the publication of peer-reviewed articles (at least 4 articles) in high-impact journals to ensure scientific rigor and visibility.

\subsubsection{Infographics and Press Releases}

To reach non-academic audiences and ensure inclusivity, the project will produce visually engaging infographics summarizing key findings in accessible language. These materials will be tailored for adolescents, educators, and community stakeholders, promoting environmental awareness and practical actions. Press releases will be prepared for major milestones, such as the publication of significant results and the organization of events. Dissemination through media outlets and institutional channels will amplify the project's visibility and foster dialogue beyond the academic sphere.

\subsubsection{White Book and Closing Seminar}

The White Book will serve as a comprehensive legacy document, compiling the project's insights, best practices, and policy recommendations aimed at educators, decision-makers, and community stakeholders. It will synthesize research findings into actionable strategies to promote sustainability in adolescence, ensuring that the knowledge generated by the project remains accessible and impactful beyond its duration.

In addition, a Closing Seminar will be organized as a high-profile event to present the results, engage stakeholders, and foster dialogue on sustainability and youth engagement. This seminar will provide an opportunity to share lessons learned, celebrate achievements, and strengthen networks for future collaboration.

% ============================================================
\section{Monitoring}
% ============================================================

To ensure effectiveness, dissemination activities will be monitored through predefined indicators such as social media engagement, number of downloads from OSF, participation in events, and citation metrics. Regular reviews will allow adjustments to maximize impact and inclusivity.
